\chapter*{Abstract}

Determinization of finite-state automata and transducers can increase the number of
states exponentially. We study finite-state devices whose graphs contain a separator: a
set of states $X$ whose removal splits the remaining states into components $Q_1,
\ldots, Q_k$ that are not connected by transitions. For such $k$-separated devices we
introduce quasi-deterministic finite-state automata and quasi-subsequential
finite-state transducers. They keep one deterministic quasi-state per component and
synchronize the components through the separator. A word is processed in $\mathcal{O}(k)$
time per symbol.
\bigbreak
We show that every $k$-separated finite-state automaton has an equivalent
quasi-deterministic finite-state automaton with at most $2^{|Q_i| + |X|}$ quasi-states in
its $i$-th component, compared with up to $2^{|Q|}$ states of a deterministic
finite-state automaton. The quasi-states of a reachable state are exactly the projections
of the corresponding state of the subset construction onto the components. A direct
construction processes as many states as the subset construction. We give a
per-component construction that runs in $\mathcal{O}((|\Sigma| + 2^{|X|}) |Q| 2^{|X|}
\sum_{i=1}^k 2^{|Q_i|})$ time instead. For finite-state automata constructed from regular
expressions, we find separators of size $2$ with components of at most $\frac{2}{3}|Q|$
states in linear time. As a consequence, every regular expression can be
quasi-determinized in $\mathcal{O}(|\Sigma| |Q| 2^{\frac{2}{3}|Q|})$ time. For trimmed
functional $k$-separated finite-state transducers whose function has bounded variation,
we construct an equivalent quasi-subsequential finite-state transducer.
\bigbreak
The main results are formalized and machine-checked in the Lean theorem prover, and the
constructions for automata are implemented and evaluated experimentally.
